% Encabezado: aquí se cargan los paquetes y se hacen definiciones adicionales
%%%%%%%%%%%%%%%%%%%%%%%%%%%%%%%%%%%%%%%%%%%%%%%%%%%%%%%%%%%%%%%%%%%

\documentclass[11pt,a4paper]{article}
\usepackage{fullpage}                          % Márgenes amplios al estilo "fullpage"
\usepackage[spanish,es-tabla]{babel}           % Tipografía y traducciones en español
\usepackage[utf8]{inputenc}                    % Permite escribir acentos y ñ directamente
\usepackage[T1]{fontenc}                       % Codificación de fuentes (mejor división silábica)
\usepackage[onehalfspacing]{setspace}          % Interlineado 1.5
\usepackage{graphicx}                          % Inclusión de gráficos
\usepackage[breaklinks=true,colorlinks=true,linkcolor=blue,urlcolor=blue,citecolor=blue]{hyperref} % Referencias e hiperenlaces
\usepackage{amsmath,amsthm,amssymb}            % Entornos y símbolos matemáticos
\usepackage{icomma}                            % Coma decimal inteligente
\usepackage[output-decimal-marker={,},per-mode=symbol]{siunitx}% Unidades del SI con coma decimal
\usepackage{placeins}                          % Comando \FloatBarrier
\usepackage[natbib,abbreviate=true,doi=false,style=numeric-comp,giveninits=true,sorting=none]{biblatex} % Bibliografía moderna (requiere biber)
\usepackage{tikz}
\usetikzlibrary{decorations.pathreplacing,calligraphy}

\addbibresource{MyBibliography.bib}            % Archivo .bib con las referencias

\graphicspath{{Figuras/}}                      % Carpeta donde se buscan las figuras

\DeclareSIUnit{\dBm}{dBm}
\DeclareSIUnit[per-mode=reciprocal]\WN{\per\centi\meter}



\title{
       \null \vspace{-35mm}
\includegraphics[width=3cm]{images/IFlogo.png}\hfill\includegraphics[width=3cm]{images/UNAMlogo.png}
\\ 
\vspace{1.0 cm}TEMAS SELECTOS DE CÓMPUTO DE ALTO RENDIMIENTO
PARA SIMULACIONES FÍSICAS\\[1ex]
\large Reporte de Simulación de Calor en 1D y 2D}
\author{Fís. Saúl Ruano Sánchez \thanks{\href{mailto:s.ruano@estudiantes.fisica.unam.mx}{s.ruano@estudiantes.fisica.unam.mx}}}
\date{Semestre 2026-2}

%%%%%%%%%%%%%%%%%%%%%%%%%%%%%%%%%%%%%%%%%%%%%%%%%%%%%%%%%%%%%%%%%%%
\begin{document}
%}


\maketitle
\vfill

\renewcommand{\abstractname}{Resumen}
\section*{\abstractname}
\textit{.. ALGO}
\thispagestyle{empty}
%
%
\tableofcontents
\thispagestyle{empty}
\cleardoublepage
\pagenumbering{arabic}
\newpage
%
%
\section{Introducción}
\label{sec:Introduccion}
%
\textit{¿De qué trata el trabajo? ¿Cuál es la pregunta principal y por qué es interesante?}\\
%
\section{Fundamentos y teoría}
\label{sec:Fundamentos}
%

\subsection{Ecuación de Calor}
El flujo de calor por difusión en un medio continuo se rige por la ley de Fourier,
que establece que el flujo de calor $\vec{J}_{q}$ es proporcinal al gradiente negativo
de la temperatura, 
\begin{equation}
       \vec{J_q} = - \kappa \nabla T,
       \label{eq:Fourier}
\end{equation}
donde $\kappa$ representa la conductividad térmica del material.
Al combinar \eqref{eq:Fourier}, con la ecuación de continuidad de la energía interna, 
en ausencias de fuentes, obtenemos la ecuación de calor 
\begin{equation}
\rho \, \, \partial_t U =  \rho \, c_p  \, \partial T =
- \nabla \cdot \left( \vec{J}_{q} \right) = k \nabla^2 T 
\quad \Rightarrow \quad 
\partial_t T = \alpha \nabla^2 T,
\label{eq:calor}
\end{equation}
con $\alpha = k/(\rho \, c_p)$. Nos centramos en el régimen estacionario para una y dos dimensiones sobre un 
dominio rectangular $\Omega$, con distintitas condiciones en la frontera $\partial \Omega$,
En este caso, la ecuación \eqref{eq:calor} toma la forma de una ecuación de Laplace,

\begin{equation}
       \nabla^2 T = 0
       \label{eq:Laplace}
\end{equation}

\subsection{Condiciones de Frontera}

Para la ecuación de Laplace, consideramos dos tipos de condiciones sobre la frontera $\partial \Omega$ del dominio rectangular:

\textbf{Condiciones tipo Dirichlet:} se impone el valor de la temperatura sobre la frontera,

\begin{equation}
T(\mathbf{x}) = T_D(\mathbf{x}), \qquad \mathbf{x} \in \partial\Omega_D,
\label{eq:Dirichlet}
\end{equation}

donde $T_D$ es una función conocida, para nuestro caso la tomaremos como constante, y $\partial \Omega_D$
denota la porción de la frontera donde se aplica la condición.


\textbf{Condiciones tipo Neumann:} se impone el valor del flujo normal de calor sobre la frontera, 

\begin{equation}
-k\frac{\partial T}{\partial n} = q_N(\mathbf{x}), \qquad \mathbf{x} \in \partial\Omega_N,
\label{eq:Neumann}
\end{equation}
o bien, de forma equivalente, la derivada normal de la temperatura:
\begin{equation}
\frac{\partial T}{\partial n} = g_N(\mathbf{x}), \qquad \mathbf{x} \in \partial\Omega_N,
\end{equation}
donde $\partial/\partial n$ denota la derivada en la dirección normal exterior a la frontera, 
$q_{N}$ es el flujo dado.

%
\subsection{Diferencias Finitas}
Para resolver numéricamente la ecuación es necesario discretizar su dominio. 
El método de diferencias finitas logra esto aproximando las derivadas mediante 
desarrollos en serie de Taylor de la función que se quiere resolver.

\subsubsection{Caso 1D}

Partiendo de la ecuación~\eqref{eq:Laplace}, que en una dimensión tiene la siguiente forma, 

\begin{equation}
       \frac{d^2 T}{dx^2} = 0,
\end{equation}       
para aproximar la segunda derivada, recurrimos al desarrollo en serie de Taylor de la función $T$
en torno a $x$, donde $\Delta x$ es el espaciamiento de la malla sobre la que se resuelve la ecuación.
En la figura~\ref{fig:malla1D} se muestra la discretización resultante.


\begin{equation}
       \frac{d^2 T}{d x^2} \approx \frac{T_{i+1} - 2 \, T_{i} + T_{i-1}}{(\Delta x)^2}
       = 0.
\end{equation}
donde adoptamos la notación,
$T_i \equiv T(x_i)$, $T_{i\pm 1} \equiv T(x_i \pm \Delta x)$. Se observa que el problema continuo 
se reduce a uno lineal
\begin{equation}
    T_{i+1} - 2\, T_i + T_{i-1} = 0, \qquad i = 1, \dots, N-2,
    \label{eq:laplaceDisc}
\end{equation}
donde $N$ es el número total de nodos en la malla. Las ecuaciones en los nodos de frontera 
$(i=0)$ y $(i=N-1)$ se sustituyen por las condiciones de fronter correspondientes.

\begin{figure}[h]
       \centering
       \includegraphics[width=\linewidth]{images/malla1D.png}
       \caption{Discretización del espacio.}
       \label{fig:malla1D}
\end{figure}

Suponiendo que queremos resolver la ecuación de Laplace en el intervalo
$[0, L]$, con condiciones de frontera de Dirichlet, de tal forma que fijamos el valor de la temperatura
en los extremos,

\begin{equation}
       T(0) = T_0, \qquad T(L) = T_L,
\end{equation}
para ejemplificar el método, discretizamos el dominio en $N=6$ nodos, con $\Delta x = L /(N-1)$.
La ecuación discretizada~\eqref{eq:laplaceDisc} se aplica en los nodos interiores $i=1,2,3,4$, 
mientras que en los nodos de frontera imponemos directamente los valores dados por las condiciones de frontera. De tal forma que obtenemos, 
el siguiente sistema de ecuaciones,
\begin{equation}
       \begin{pmatrix}
              1 & 0 & 0 & 0 & 0 & 0 \\
              1 & -2 & 1 & 0 & 0 & 0 \\
              0 & 1 & -2 & 1 & 0 & 0 \\
              0 & 0 & 1 & -2 & 1 & 0 \\
              0 & 0 & 0 & 1 & -2 & 1 \\
              0 & 0 & 0 & 0 & 0  & 1 
       \end{pmatrix}
       \begin{pmatrix}
              T_0 \\
              T_1 \\
              T_2 \\ 
              T_3 \\
              T_4 \\
              T_5 \\
       \end{pmatrix}
       = 
       \begin{pmatrix}
              T_0 \\
              0 \\
              0 \\ 
              0 \\
              0 \\
              T_L \\
       \end{pmatrix}
\end{equation}
Suponiendo ahora condiciones frontera de Neumann, en los dos extremos,
\begin{equation}
\frac{dT}{dx}\bigg|_{x=0} = g_0, \qquad \frac{dT}{dx}\bigg|_{x=L} = g_L,
\end{equation}
el sistema ecuaciones queda prácticamente igual salvo por las ecuaciones asociadas a los nodos frontera. 
En cada extremo, se aproxima la derivada mediante una diferencia adelanta en $x=0$
\begin{equation*}
       \frac{dT}{dx} \approx \frac{T_1 - T_0}{\Delta x} = g_0 
       \quad \Rightarrow \quad
       T_1 - T_0 = (g_0)\Delta x,
\end{equation*}
y mediante una diferencia atrasada en $x=L$
\begin{equation*}
       \frac{dT}{dx} \approx \frac{T_5 - T_4}{\Delta x} = g_0 
       \quad \Rightarrow \quad
       T_5 - T_4 = (g_L)\Delta x.
\end{equation*}
Se obtiene el siguiente sistema de ecuaciones,
\begin{equation}
       \begin{pmatrix}
              -1 & 1 & 0 & 0 & 0 & 0 \\
              1 & -2 & 1 & 0 & 0 & 0 \\
              0 & 1 & -2 & 1 & 0 & 0 \\
              0 & 0 & 1 & -2 & 1 & 0 \\
              0 & 0 & 0 & 1 & -2 & 1 \\
              0 & 0 & 0 & 0 & -1  & 1 
       \end{pmatrix}
       \begin{pmatrix}
              T_0 \\
              T_1 \\
              T_2 \\ 
              T_3 \\
              T_4 \\
              T_5 \\
       \end{pmatrix}
       = 
       \begin{pmatrix}
              (g_0) \Delta x \\
              0 \\
              0 \\ 
              0 \\
              0 \\
              (g_L) \Delta x \\
       \end{pmatrix},
\end{equation}
para que la solución tenga sentido físicamente tenemos que la suma de los flujos en la fronteras sea igual a cero, $g_0 + g_L = 0$.
lo cual expresa que el flujo de calor que entra al sistema es el mismo que sale. Dicho de otra forma, en estado estacionario y sin fuentes internas, 
la energía que ingresa por una frontera debe ser evacuada en su totalidad por la otra; de lo contrario, el sistema no alcanzaría el equilibrio térmico y 
la temperatura crecería o decrecería indefinidamente.

\subsubsection{Caso 2D}



\section{Metodología e implementación}
\label{sec:Metodologia}
%
\textit{¿Qué se necesitaría para reproducir el trabajo? Describa algoritmos, hardware, bibliotecas y configuración.}\\
Como ejemplo de ecuación, considere la siguiente expresión~\cite{TiltFilter}:
\begin{equation}
\lambda(\Phi)=\lambda(0)\sqrt{1-\frac{\sin^2\Phi}{n_\text{ef}^2}}\;.
\label{eq:ejemplo}
\end{equation}
donde $n_\text{ef}$ es el índice de refracción efectivo.
%
\section{Resultados}
\label{sec:Resultados}
%
\textit{¿Cuáles son las respuestas medidas a la(s) pregunta(s) principal(es)?}\\
% La Figura~\ref{fig:ejemplo} muestra un ejemplo.
% \begin{figure}[htb!]
%  \centering
%  \includegraphics[width=0.65\textwidth]{ejemplo_grafica}
%  \caption{\label{fig:ejemplo} Descripción de la figura.}
% \end{figure}
%
\section{Conclusiones}
\label{sec:Conclusiones}
%
\textit{¿Cuál es la respuesta final? ¿Cómo podría mejorarse el trabajo?}\\
La referencia~\cite{GP1StromSpannung} es un ejemplo de cómo citar un manual.
%

\printbibliography
\vfill

\section*{Declaración}

Por medio de la presente declaro que este reporte fue elaborado de forma independiente y que se han indicado todas las fuentes y referencias necesarias.

\begin{tabular}{@{}p{2.5in}p{2.5in}@{}}
 \\[5\bigskipamount]
  \dotfill & \dotfill \\
  Estudiante & Fecha \\
\end{tabular}

\end{document}
